% AUTHORED
\section{Gradient descent, SGD, Adam}\label{sec:06}
\conceptbadge{06}

Gradients tell you a local direction in parameter space. Optimization decides
how far to move. For parameters $\theta$ and loss $L$, vanilla gradient descent
updates
\[
  \theta \leftarrow \theta - \eta \nabla_\theta L,
\]
where $\eta$ is the learning rate. The minus sign matters: the gradient points
toward the steepest local increase, so training moves in the opposite direction.

\subsection*{Mini-batches}
The exact gradient over a full dataset can be expensive. Stochastic gradient
descent (SGD) estimates it from one example or a mini-batch. A mini-batch
gradient is noisy, but it is much cheaper, and the noise can help optimization
move out of poor local minima. In language modeling, each batch contains many
short token sequences, and each position contributes a next-token loss.

\subsection*{Learning rate}
The learning rate is the most important optimizer hyperparameter. If it is
too small, training is very slow. If it is too large, the loss can oscillate or
diverge. Practical
training often uses schedules: warm up early steps, decay later steps, or lower
the rate when validation loss stops improving.

\subsection*{Momentum and Adam}
Momentum keeps an exponential moving average of recent gradients and uses it as
velocity. Adam keeps two moving averages: one for gradients and one for squared
gradients. The second moment rescales updates so parameters with consistently
large gradients take smaller effective steps. Adam is a common default for
transformers because it is robust across many parameter groups and scales.

\subsection*{Training loop discipline}
A minimal loop has the same shape throughout the course:
\begin{enumerate}
  \item sample a mini-batch;
  \item run the forward pass;
  \item compute the loss;
  \item clear old gradients;
  \item run backward;
  \item update parameters.
\end{enumerate}
Gradients accumulate by design in autograd systems, so clearing them is not
optional. Evaluation uses the same forward pass but no parameter update.

\subsection*{Exercise connection}
The exercise isolates the simplest update, \texttt{params - lr * grads}. P2
will expand that idea into optimizer classes that update all trainable
parameters in the project.

\medskip
\noindent\textit{Go deeper:} Goodfellow, Bengio \& Courville, \emph{Deep
Learning}, Ch.~8 (\texttt{https://www.deeplearningbook.org}) covers optimization
for deep models, including momentum and Adam.
